\section{研究议程：下一代多模态 Agent 的三个突破口}

\subsection{突破口一：统一表示层}
课程大量讨论了截图、可访问性树、动作代码和推理链，这些本质上都是不同模态的中间表示。下一阶段的重要方向是建立统一表示层，减少模态切换的信息损耗。

\begin{knowledgebox}{统一表示层要解决的问题}
\begin{itemize}
    \item 如何把视觉对象、语义对象、动作对象映射到同一结构；
    \item 如何让不同平台（web、desktop、mobile）共享抽象接口；
    \item 如何支持长链路任务中的状态压缩与可逆恢复。
\end{itemize}
\end{knowledgebox}

\subsection{突破口二：可验证长程规划}
现有模型在复杂任务上仍容易陷入短视策略。课程后段强调 planner 贡献，说明 ``可验证规划'' 会成为下一轮竞争焦点：不是只生成计划，而是持续验证计划是否仍与环境状态一致。

\begin{importantbox}{规划能力的评测应从 ``有无'' 进化到 ``质量''}
未来评测不仅要问模型能否给出 plan，还要评估 plan 的可执行性、可更新性和错误恢复效率。没有这些指标，planner 容易沦为装饰性中间文本。
\end{importantbox}

\subsection{突破口三：安全与治理原生化}
当 Agent 开始操作真实账户和业务系统，安全不可能后置。权限边界、审计日志、异常回滚和人类接管机制必须成为模型训练和部署时的原生约束，而非上线前补丁。

\begin{warningbox}{Agent 的最大风险不只是 ``做不到''，而是 ``做错且不可追溯''}
如果系统缺少可审计链路，事故发生后无法复盘责任边界。课程虽然聚焦研究，但其环境与评估思想已经为治理原生化提供了实践基础。
\end{warningbox}

\subsection{未来 12 个月可执行任务清单}
\begin{table}[H]
\centering
\begin{tabular}{p{2.8cm}p{4.8cm}p{3.6cm}}
\toprule
\textbf{方向} & \textbf{建议任务} & \textbf{预期产出} \\
\midrule
统一表示 & 构建跨平台对象-动作中间层 & 降低迁移成本与动作错误率 \\
规划优化 & 引入在线计划校验与重规划机制 & 提升长任务完成率 \\
数据飞轮 & 失败样本自动分桶并回流训练 & 降低高频失败复现率 \\
安全治理 & 建立权限模板、审计与回滚策略 & 支撑生产可控部署 \\
\bottomrule
\end{tabular}
\caption{从课程方法论落到团队执行层的路线图}
\end{table}

\subsection{本章小结}
下一代多模态 Agent 的竞争，不会只看模型参数规模，而会看谁能把统一表示、长程规划和治理机制做成可复用工程能力。

